This section creates the brain of the system: one Lambda function
handling all business logic, and an API Gateway as the entry point.

\textbf{Security Note (Auth):} In the scope of this workshop, password
hashing and user management logic are coded directly within the API so
you can easily visualize the data flow. In a real-world Production
environment, you should delegate this entire functionality to
\textbf{Amazon Cognito} to meet security standards.

Order matters: \textbf{IAM role first, then Lambda, then API Gateway}.
Because Lambda needs the role at creation time, and API Gateway needs
the Lambda to already exist.

\vspace{0.5em}\noindent\textbf{Step 1. Create an IAM role with least
privilege}\par

Lambda needs read/write access to DynamoDB and permission to call
Personalize. The quickest option would be attaching
\texttt{AdministratorAccess}, but that is a bad habit: if your code has
a vulnerability, an attacker gains the whole account.

We grant only what is needed.

\begin{enumerate}
\def\labelenumi{\arabic{enumi}.}
\tightlist
\item
  Open the \href{https://console.aws.amazon.com/iam/home\#/roles}{IAM
  console}, choose \textbf{Roles} then \textbf{Create role}
\item
  \textbf{Trusted entity type}: AWS service. \textbf{Use case}: Lambda.
  Choose \textbf{Next}
\item
  Search for and select \texttt{AWSLambdaBasicExecutionRole} ---
  permission to write CloudWatch logs
\item
  Choose \textbf{Next}, name it \texttt{fcj-recsys-lambda-role}, choose
  \textbf{Create role}
\end{enumerate}

Now add the specific permissions for DynamoDB and Personalize:

\begin{enumerate}
\def\labelenumi{\arabic{enumi}.}
\setcounter{enumi}{4}
\tightlist
\item
  Open the new role, choose \textbf{Add permissions} then \textbf{Create
  inline policy}
\item
  Switch to the \textbf{JSON} tab and paste the following (replace
  \texttt{\textless{}ACCOUNT-ID\textgreater{}} with your account
  number):
\end{enumerate}

\begin{Shaded}
\begin{Highlighting}[]
\FunctionTok{\{}
  \DataTypeTok{"Version"}\FunctionTok{:} \StringTok{"2012{-}10{-}17"}\FunctionTok{,}
  \DataTypeTok{"Statement"}\FunctionTok{:} \OtherTok{[}
    \FunctionTok{\{}
      \DataTypeTok{"Effect"}\FunctionTok{:} \StringTok{"Allow"}\FunctionTok{,}
      \DataTypeTok{"Action"}\FunctionTok{:} \OtherTok{[}
        \StringTok{"dynamodb:GetItem"}\OtherTok{,} \StringTok{"dynamodb:PutItem"}\OtherTok{,} \StringTok{"dynamodb:UpdateItem"}\OtherTok{,}
        \StringTok{"dynamodb:DeleteItem"}\OtherTok{,} \StringTok{"dynamodb:Query"}\OtherTok{,} \StringTok{"dynamodb:Scan"}\OtherTok{,}
        \StringTok{"dynamodb:BatchGetItem"}\OtherTok{,} \StringTok{"dynamodb:BatchWriteItem"}
      \OtherTok{]}\FunctionTok{,}
      \DataTypeTok{"Resource"}\FunctionTok{:} \StringTok{"arn:aws:dynamodb:ap{-}southeast{-}1:\textless{}ACCOUNT{-}ID\textgreater{}:table/*"}
    \FunctionTok{\}}\OtherTok{,}
    \FunctionTok{\{}
      \DataTypeTok{"Effect"}\FunctionTok{:} \StringTok{"Allow"}\FunctionTok{,}
      \DataTypeTok{"Action"}\FunctionTok{:} \StringTok{"personalize:GetRecommendations"}\FunctionTok{,}
      \DataTypeTok{"Resource"}\FunctionTok{:} \StringTok{"*"}
    \FunctionTok{\}}
  \OtherTok{]}
\FunctionTok{\}}
\end{Highlighting}
\end{Shaded}

\begin{enumerate}
\def\labelenumi{\arabic{enumi}.}
\setcounter{enumi}{6}
\tightlist
\item
  Name the policy \texttt{fcj-recsys-lambda-policy} and choose
  \textbf{Create policy}
\end{enumerate}

Get your Account ID with
\texttt{aws\ sts\ get-caller-identity\ -\/-query\ Account\ -\/-output\ text}

\vspace{0.5em}\noindent\textbf{Step 2. Package the source
code}\par

Lambda needs a zip file containing your code and its dependencies:

\begin{Shaded}
\begin{Highlighting}[]
\BuiltInTok{cd}\NormalTok{ backend}
\ExtensionTok{npm}\NormalTok{ install }\AttributeTok{{-}{-}production}
\FunctionTok{zip} \AttributeTok{{-}r}\NormalTok{ ../function.zip index.mjs node\_modules}
\end{Highlighting}
\end{Shaded}

\vspace{0.5em}\noindent\textbf{Step 3. Create the Lambda
function}\par

\begin{enumerate}
\def\labelenumi{\arabic{enumi}.}
\tightlist
\item
  Open the
  \href{https://ap-southeast-1.console.aws.amazon.com/lambda/home?region=ap-southeast-1}{Lambda
  console}
\item
  Choose \textbf{Create function}
\item
  Fill in:

  \begin{itemize}
  \tightlist
  \item
    \textbf{Function name}: \texttt{fcj-recsys-api}
  \item
    \textbf{Runtime}: \textbf{Node.js 20.x}
  \item
    \textbf{Architecture}: x86\_64
  \item
    \textbf{Execution role}: choose \textbf{Use an existing role},
    select \texttt{fcj-recsys-lambda-role}
  \end{itemize}
\item
  Choose \textbf{Create function}
\item
  On the \textbf{Code} tab, choose \textbf{Upload from} then
  \textbf{.zip file}, and upload \texttt{function.zip}
\item
  Go to \textbf{Runtime settings}, choose \textbf{Edit}, set
  \textbf{Handler} to \texttt{index.handler}
\end{enumerate}

\begin{figure}
\centering
\pandocbounded{\includegraphics[keepaspectratio,alt={Create Lambda}]{5-Workshop/5.5/create-lambda.png}}
\caption{Create Lambda}
\end{figure}

\vspace{0.5em}\noindent\textbf{Step 4. Increase the
timeout}\par

By default Lambda stops after 3 seconds, which is not enough when
calling Personalize.

\begin{enumerate}
\def\labelenumi{\arabic{enumi}.}
\tightlist
\item
  Open the \textbf{Configuration} tab, choose \textbf{General
  configuration}, choose \textbf{Edit}
\item
  Set \textbf{Timeout} to \texttt{30} seconds and \textbf{Memory} to
  \texttt{512} MB
\item
  Choose \textbf{Save}
\end{enumerate}

\vspace{0.5em}\noindent\textbf{Step 5. Create the HTTP
API}\par

\begin{enumerate}
\def\labelenumi{\arabic{enumi}.}
\tightlist
\item
  Open the
  \href{https://ap-southeast-1.console.aws.amazon.com/apigateway/main/apis?region=ap-southeast-1}{API
  Gateway console}
\item
  In the \textbf{HTTP API} box, choose \textbf{Build}
\item
  Choose \textbf{Add integration}, type \textbf{Lambda}, and select the
  \texttt{fcj-recsys-api} function
\item
  \textbf{API name}: \texttt{fcj-recsys-api}, choose \textbf{Next}
\item
  Under \textbf{Configure routes}, define a catch-all route:

  \begin{itemize}
  \tightlist
  \item
    \textbf{Method}: \texttt{ANY}
  \item
    \textbf{Resource path}: \texttt{/\{proxy+\}}
  \item
    \textbf{Integration target}: the Lambda function you selected
  \end{itemize}
\item
  Choose \textbf{Next}, keep the \texttt{\$default} stage with
  \textbf{Auto-deploy} enabled, choose \textbf{Create}
\end{enumerate}

The \texttt{/\{proxy+\}} route sends every path to the same Lambda,
letting the code route internally. This keeps API Gateway configuration
simple, at the cost of moving routing into the code. With 22 routes,
that is a reasonable trade-off.

\vspace{0.5em}\noindent\textbf{Step 6. Enable CORS}\par

The frontend runs on a CloudFront domain while the API sits on a
different one. Browsers will block the calls unless CORS is declared.

\begin{enumerate}
\def\labelenumi{\arabic{enumi}.}
\tightlist
\item
  In your API, choose \textbf{CORS} in the left menu, then
  \textbf{Configure}
\item
  Set:

  \begin{itemize}
  \tightlist
  \item
    \textbf{Access-Control-Allow-Origin}: \texttt{*}
  \item
    \textbf{Access-Control-Allow-Headers}:
    \texttt{content-type,\ authorization}
  \item
    \textbf{Access-Control-Allow-Methods}:
    \texttt{GET,\ POST,\ PUT,\ DELETE,\ OPTIONS}
  \end{itemize}
\item
  Choose \textbf{Save}
\end{enumerate}

CORS errors only appear in the browser console; testing with
\texttt{curl} still works fine. If the interface reports it cannot reach
the server while \texttt{curl} succeeds, press F12 and check the Console
tab first.

\vspace{0.5em}\noindent\textbf{Step 7. Connect the frontend to the
API}\par

Copy the \textbf{Invoke URL} from the API detail page, then go back to
the frontend folder:

\begin{Shaded}
\begin{Highlighting}[]
\BuiltInTok{cd}\NormalTok{ frontend}
\BuiltInTok{echo} \StringTok{"VITE\_API\_BASE\_URL=https://xxxxx.execute{-}api.ap{-}southeast{-}1.amazonaws.com"} \OperatorTok{\textgreater{}}\NormalTok{ .env}
\ExtensionTok{npm}\NormalTok{ run build}
\ExtensionTok{aws}\NormalTok{ s3 sync dist/ s3://fcj{-}recsys{-}frontend{-}}\OperatorTok{\textless{}}\NormalTok{your{-}name}\OperatorTok{\textgreater{}}\NormalTok{/ }\AttributeTok{{-}{-}delete}
\ExtensionTok{aws}\NormalTok{ cloudfront create{-}invalidation }\AttributeTok{{-}{-}distribution{-}id} \OperatorTok{\textless{}}\NormalTok{YOUR{-}ID}\OperatorTok{\textgreater{}}\NormalTok{ {-}{-}paths }\StringTok{"/*"}
\end{Highlighting}
\end{Shaded}

\vspace{0.5em}\noindent\textbf{Verify}\par

\begin{Shaded}
\begin{Highlighting}[]
\ExtensionTok{curl}\NormalTok{ https://xxxxx.execute{-}api.ap{-}southeast{-}1.amazonaws.com/products}
\end{Highlighting}
\end{Shaded}

A JSON list of products means success. Reload the CloudFront page and
products should now appear.
